\documentclass{article}
\usepackage{graphicx}
\usepackage[T1]{fontenc}
\usepackage[utf8]{inputenc}
\usepackage[spanish]{babel}
\usepackage{titling}
\setlength{\droptitle}{-2cm}
\usepackage[a4paper,top=3.0cm,bottom=2.5cm,left=2.5cm,right=2.5cm]{geometry}

%% Ref
\usepackage{hyperref}
\hypersetup{
  colorlinks=true,
  linkcolor=blue,
  urlcolor=blue,
  citecolor=blue
}

%% Profundidad de la tabla de contenidos:
%% 1 = \section, 2 = \subsection, 3 = \subsubsection
\setcounter{tocdepth}{2}

%% Title
\title{Contenido del curso}
\author{Analítica para los negocios · 06327-ECO}
\date{Periodo 2026-02}

%%=======================================%%
\begin{document}
\maketitle
\tableofcontents
\clearpage

%% Nota metodológica (aplica a todas las semanas de contenido)
\noindent\textbf{Dinámica semanal.} En las semanas de contenido, la fundamentación teórica se estudia \textbf{antes} de clase mediante un podcast o un video (los enlaces estarán disponibles en Intu) y se evalúa con un \textbf{quiz} semanal. La sesión de clase se dedica a la \textbf{aplicación} de esa teoría (aplicación guiada / taller).
\vspace{0.3cm}

%%=======================================%%
%%=======================================%%
\section{Fundamentos del curso}

%%=======================================%%
\subsection{Semana 01 - Presentación del curso y metodología}
\vspace{0.2 cm}

\begin{itemize}
\item \textbf{1. Presentación del profesor}
\begin{itemize}
  \item Perfil y experiencia relevante para el curso.
\end{itemize}

\item \textbf{2. Presentación de estudiantes}
\begin{itemize}
  \item Ronda rápida: nombre, programa, ¿hace cuánto tiempo vio BA-2?
  \item ¿Ha tenido experiencia con datos/R?
\end{itemize}

\item \textbf{3. Motivación}
\begin{itemize}
  \item ¿Por qué Business Analytics?
  \item ¿Qué problemas ayuda a resolver en economía/negocios?
\end{itemize}

\item \textbf{4. ¿Qué es Business Analytics? (marco del curso)}
\begin{itemize}
  \item Definición operativa.
  \item Proceso analítico (alto nivel).
  \item Tipos de analítica (alto nivel).
\end{itemize}

\item \textbf{5. Contenido del curso (vista general semana a semana)}
\begin{itemize}
  \item Recorrido por unidades y semanas.
  \item Conexión con el proyecto final.
\end{itemize}

\item \textbf{6. Dinámica de aprendizaje (cómo vamos a aprender)}
\begin{itemize}
  \item Teoría antes de clase: cada semana de contenido tiene un podcast o video con la fundamentación teórica.
  \item Quiz semanal sobre esa teoría (en Intu).
  \item La sesión de clase se dedica a la aplicación (aplicación guiada / talleres).
  \item Trabajo en clase y trabajo autónomo; expectativas de participación.
\end{itemize}

\item \textbf{7. Evaluación y reglas de juego}
\begin{itemize}
  \item Componentes: quizzes semanales, talleres en clase, dos exámenes parciales y proyecto final (tres entregas).
  \item Todas las semanas de contenido hay quiz sobre el material teórico.
  \item No hay recuperación de quizzes ni de talleres; únicamente los exámenes parciales tienen supletorio (según reglamento).
  \item Se elimina la peor nota de quizzes y la peor nota de talleres \textbf{una sola vez} en el semestre.
  \item Entregas del proyecto: quien no presenta una entrega obtiene cero.
\end{itemize}

\item \textbf{8. Trabajo final}
\begin{itemize}
  \item Datos: una base de datos empresarial con múltiples variables; cada grupo define su propia pregunta de negocio (clasificación, regresión o segmentación).
  \item Entregables: (1) pregunta de negocio; (2) presentación del avance/EDA (Semana 09); (3) presentación final (Semana 16).
  \item La entrega final es una sustentación oral (ya no hay documento escrito) y equivale a un examen parcial: la ausencia solo se acepta mediante supletorio autorizado por la universidad.
  \item Contenidos mínimos, rúbricas y reglas: documento de lineamientos del proyecto.
\end{itemize}

\item \textbf{9. Política de IA}
\begin{itemize}
  \item Qué se permite / qué no se permite.
  \item Cómo citar/explicar uso de IA.
  \item Enfoque en aprendizaje y trazabilidad.
\end{itemize}
\end{itemize}

%%=======================================%%
\newpage
\subsection{Semana 02 - Introducción a LLMs}
\vspace{0.2 cm}

\begin{itemize}
\item \textbf{Alcance}
\begin{itemize}
  \item Qué es un LLM (\textit{Large Language Model}) y por qué es relevante para economía y negocios.
  \item La IA como herramienta transversal del curso: apoyo para R, visualización, proceso analítico y proyecto final.
  \item Pasar de instrucciones vagas (``hágame esto'') a prompts efectivos que le sacan provecho a la IA.
\end{itemize}

\item \textbf{Fundamentación Teórica (podcast + quiz)}
\begin{itemize}
  \item Podcast: \emph{¿Qué es un LLM?} % TODO: enlace al podcast en Intu
  \begin{itemize}
    \item Intuición de cómo funciona un LLM y su arquitectura (alto nivel).
    \item Alcances y límites: qué hace bien, dónde se equivoca (alucinaciones) y por qué hay que verificar.
    \item Buenas prácticas de uso en contextos académicos y de negocio.
  \end{itemize}
  \item Material para la sesión: % TODO: pegar aquí el/los enlace(s) de podcast/video de esta semana
  \item Quiz de teoría del podcast (en Intu).
\end{itemize}

\item \textbf{Aplicación en clase}
\begin{itemize}
  \item Cómo escribir un buen prompt: contexto, tarea, formato de salida y restricciones; iterar sobre la respuesta.
  \item Introducción a \textit{skills} (instrucciones reutilizables) para tareas recurrentes.
  \item Ejercicios de prompts aplicados a tareas del curso (pedir ayuda en R, en visualización, en el proceso analítico).
\end{itemize}
\end{itemize}

%%=======================================%%
\newpage
\subsection{Semana 03 - Fundamentos de R}
\vspace{0.2 cm}

\begin{itemize}
\item \textbf{Alcance}
\begin{itemize}
  \item Familiarizarse con la interfaz de RStudio: consola, editor (scripts) y paneles principales.
  \item Usar R como calculadora: operaciones aritméticas, comparaciones y lógica básica.
  \item Reconocer tipos de datos (numérico, texto, lógico) y valores especiales (NA/NULL).
  \item Entender qué es un \emph{objeto} en R y cómo asignarlo (operador \texttt{<-}).
  \item Trabajar con estructuras básicas: \textbf{vectores}, \textbf{matrices} (intro) y \textbf{data frames}.
  \item Introducción mínima a funciones, sistema de ayuda (\texttt{?}, \texttt{help()}) y carga de librerías.
  \item Cierre: noción de entorno de trabajo (Environment) y buenas prácticas iniciales (Projects / rutas relativas).
\end{itemize}

\item \textbf{Fundamentación Teórica (video + quiz)}
\begin{itemize}
  \item Video de clase: recorrido por la interfaz y los conceptos base. % TODO: enlace al video en Intu
  \begin{itemize}
    \item Interfaz de trabajo en RStudio y flujo recomendado: \emph{script} $\rightarrow$ ejecución $\rightarrow$ revisión en consola/environment.
    \item Tipos de datos y coerción básica (numérico, carácter, lógico; NA/NULL).
    \item Objetos en R: asignación, nombres, y funciones para inspección (\texttt{class()}, \texttt{typeof()}, \texttt{str()}).
    \item Estructuras de datos: vectores (creación e indexación desde 1), matrices (idea general) y data frames (acceso básico a filas/columnas).
    \item Paquetes: diferencia entre instalar (\texttt{install.packages}) y cargar (\texttt{library}/\texttt{require}); atajo opcional con \texttt{pacman::p\_load()}.
  \end{itemize}
  \item Material para la sesión: % TODO: pegar aquí el/los enlace(s) de podcast/video de esta semana
  \item Quiz de teoría (en Intu).
\end{itemize}

\item \textbf{Aplicación en clase (aplicación guiada)}
\begin{itemize}
  \item Mini-lab guiado (en parejas):
  \begin{itemize}
    \item Ejecutar operaciones en consola y luego repetirlas desde un \texttt{script}.
    \item Crear objetos (\texttt{x <- ...}), inspeccionar con \texttt{class/typeof/str}.
    \item Construir un vector y practicar indexación con \texttt{[]}.
    \item Crear un data frame pequeño (3 columnas) y realizar accesos básicos (\texttt{\$} y \texttt{[ , ]}).
    \item Provocar 1 \emph{warning} y 1 \emph{error} y aprender a interpretarlos.
    \item Cargar una librería (o \texttt{pacman::p\_load}) y verificar que quedó activa.
  \end{itemize}
  \item Cierre (10 min): conectar Environment con los objetos creados y explicar por qué usaremos Projects/rutas relativas para reproducibilidad.
\end{itemize}
\end{itemize}

%%=======================================%%
\newpage
\subsection{Semana 04 - Manipulación y visualización: \texttt{dplyr} + \texttt{ggplot2}}
\vspace{0.2 cm}

\begin{itemize}
\item \textbf{Alcance}
\begin{itemize}
  \item Manipulación de \textbf{data frames limpios} (sin NA, sin outliers, sin inconsistencias): enfoque en \textbf{transformación} y \textbf{resumen}.
  \item \texttt{dplyr} como gramática para transformar datos: filtrar filas, seleccionar columnas, ordenar, renombrar, crear y eliminar variables.
  \item Resúmenes y descriptivas: estadísticas globales y por grupo (\texttt{group\_by()} + \texttt{summarise()}), conteos con \texttt{count()} y cardinalidad con \texttt{n\_distinct()}.
  \item \texttt{ggplot2} como sistema de construcción de gráficos por capas: \texttt{ggplot()} + \texttt{aes()}, geometrías (\texttt{geom\_*()}), etiquetas (\texttt{labs()}), temas (\texttt{theme\_*()}) y escalas (\texttt{scale\_*()}).
  \item Producto de la semana: una \textbf{tabla analítica de KPIs} y 2--3 gráficos que la comunican.
  \item Buenas prácticas de comunicación: claridad en ejes, unidades, títulos y leyendas; evitar gráficos confusos o sobrecargados.
\end{itemize}

\item \textbf{Fundamentación Teórica (video + quiz)}
\begin{itemize}
  \item Video de clase. % TODO: enlace al video en Intu
  \begin{itemize}
    \item Verbo a verbo de \texttt{dplyr}: \texttt{select()}, \texttt{rename()}, \texttt{filter()}, \texttt{arrange()}, \texttt{mutate()}, \texttt{summarise()}, \texttt{group\_by()}, \texttt{count()}, \texttt{distinct()}.
    \item Resúmenes típicos: medidas de centro (media/mediana), dispersión (desviación estándar, mín, máx) y tamaño de muestra por grupo (\texttt{n()}).
    \item Gramática de gráficos: datos $\rightarrow$ mapeo estético (\texttt{aes}) $\rightarrow$ geometrías $\rightarrow$ escalas/etiquetas/tema.
    \item Construcción incremental de un gráfico: ejes y estéticos, geometría según el objetivo (\texttt{geom\_point()}, \texttt{geom\_histogram()}, \texttt{geom\_col()}, \texttt{geom\_line()}), etiquetas y estilo.
    \item Opcional (si hay tiempo): facetas para comparar grupos (\texttt{facet\_wrap()} / \texttt{facet\_grid()}).
  \end{itemize}
  \item Material para la sesión: % TODO: pegar aquí el/los enlace(s) de podcast/video de esta semana
  \item Quiz de teoría (en Intu).
\end{itemize}

\item \textbf{Aplicación en clase (aplicación guiada)}
\begin{itemize}
  \item Mini-lab guiado usando una base \textbf{limpia}:
  \begin{itemize}
    \item Transformaciones básicas: seleccionar, filtrar, renombrar, ordenar y crear 1--2 variables derivadas.
    \item KPIs globales (\texttt{summarise()} sin agrupar) y KPIs por grupo (\texttt{group\_by()} + \texttt{summarise()}, incluye \texttt{n()}).
    \item Construir 2--3 gráficos con el enfoque incremental (gráfico base $\rightarrow$ geometría $\rightarrow$ etiquetas $\rightarrow$ tema).
  \end{itemize}
  \item Producto mínimo: un script reproducible con la tabla de KPIs y 2--3 visualizaciones con títulos, ejes etiquetados y leyenda clara (cuando aplique).
  \item Cierre (5--10 min): discusión de qué KPIs son útiles para una decisión de negocio y cómo comunicarlos.
\end{itemize}
\end{itemize}

%%=======================================%%
%%=======================================%%
\newpage
\section{Proceso analítico y exploración de datos}

%%=======================================%%
\subsection{Semana 05 - Proceso analítico y tipos de analítica}
\vspace{0.2 cm}

\begin{itemize}
\item \textbf{Alcance}
\begin{itemize}
  \item \textit{Business Analytics} como un proceso que transforma datos en conocimiento accionable para mejorar decisiones.
  \item Componentes del Business Analytics:
  \begin{itemize}
    \item Estadística aplicada (patrones, relaciones, significancia).
    \item Programación y tecnología (automatización, escalabilidad, reproducibilidad).
    \item Ciencia de datos (modelos predictivos y prescriptivos).
    \item Enfoque de negocio (la analítica está al servicio de una decisión).
  \end{itemize}
  \item Tipos de preguntas de negocio y su traducción a tareas analíticas: Supervisada - Clasificación y Predicción. No Supervisada - Segmentación. Mixto - Detección de anomalías. Prescriptiva - Optimización.
  \item El flujo analítico: de la pregunta a la decisión (visión general del ciclo completo).
  \item Roles en analítica y su contribución al proyecto: Data Engineer (preparación/ETL y calidad de datos), Data Analyst (EDA, visualización, reportes), Data Scientist (modelado, validación, experimentación) y Business Analyst / Data Translator (traducción técnico--negocio y recomendaciones).
\end{itemize}

\item \textbf{Fundamentación Teórica (podcast + quiz)}
\begin{itemize}
  \item Podcast de la clase. % TODO: enlace al podcast en Intu
  \begin{itemize}
    \item ¿Qué es Business Analytics? (definición y componentes).
    \item Tareas analíticas típicas y ejemplos de preguntas (clasificación, predicción, segmentación, anomalías, optimización).
    \item Flujo analítico: formulación de la pregunta de negocio, recolección de datos, limpieza y transformación, análisis exploratorio y modelado, visualización y comunicación, decisión basada en evidencia.
    \item Fases del proceso analítico: preparación de datos (ETL/ELT), exploración (EDA), modelado y resultados (interpretación, comunicación, recomendación y monitoreo).
    \item Roles: responsabilidades, entregables típicos y cómo interactúan en un mismo proyecto.
  \end{itemize}
  \item Material para la sesión: % TODO: pegar aquí el/los enlace(s) de podcast/video de esta semana
  \item Quiz de teoría (en Intu).
\end{itemize}

\item \textbf{Aplicación en clase (presentación guiada de una pregunta de negocio)}
\begin{itemize}
  \item Recorrido completo del proceso analítico a partir de una pregunta de negocio, usando como ejemplo presentaciones de proyectos de semestres anteriores.
  \item Evaluación de preguntas candidatas: de un conjunto de preguntas, identificar cuál está bien formulada (claridad, accionabilidad, factibilidad de datos) y por qué.
  \item Cierre: cada grupo esboza su propia pregunta de negocio con apoyo de los monitores; este ejercicio es insumo directo para la Entrega 1 del proyecto final.
\end{itemize}
\end{itemize}

%%=======================================%%
\newpage
\subsection{Semana 06 - EDA: fuentes, limpieza y exploración}
\vspace{0.2 cm}

\begin{itemize}
\item \textbf{Alcance}
\begin{itemize}
  \item Introducción a \textbf{fuentes de datos} y calidad de datos en contextos reales (datos no ideales).
  \item Checklist inicial de revisión: tipos de variables, rangos plausibles, valores imposibles; valores faltantes (\texttt{NA}), categorías inconsistentes, duplicados; outliers y observaciones atípicas.
  \item Limpieza básica (sin entrar aún en modelado): estandarizar nombres y tipos, tratar valores faltantes (estrategias básicas y transparentes), identificar y manejar outliers (diagnóstico y decisiones justificadas).
  \item EDA mínima: resúmenes descriptivos y 3--5 visualizaciones para entender el dataset; hallazgos y preguntas abiertas que guían el análisis posterior.
  \item Producto reproducible: \textbf{raw $\rightarrow$ clean $\rightarrow$ analysis-ready} y registro de decisiones.
\end{itemize}

\item \textbf{Fundamentación Teórica (podcast + quiz)}
\begin{itemize}
  \item Podcast de la clase (énfasis conceptual; la parte en R se trabaja en la aplicación de clase). % TODO: enlace al podcast en Intu
  \begin{itemize}
    \item Calidad de datos: dimensiones prácticas (completitud, consistencia, validez, unicidad).
    \item Revisión estructural: dimensiones, nombres, tipos y diccionario (si existe); chequeos rápidos (conteos, rangos, cardinalidad).
    \item Valores faltantes: diagnóstico (conteos/proporciones por variable) y estrategias básicas (eliminar cuando es razonable, imputación simple, indicadores de faltante). Principio: toda decisión debe quedar documentada.
    \item Outliers: detección simple (reglas de rango, IQR/boxplot, visualización) y tratamiento (verificar si es error o caso real; recorte/capping o exclusión justificada).
    \item EDA: resúmenes globales y por grupo; gráficos exploratorios para distribuciones, comparaciones y relaciones. Conclusión: la EDA produce hipótesis y decisiones de limpieza.
  \end{itemize}
  \item Material para la sesión: % TODO: pegar aquí el/los enlace(s) de podcast/video de esta semana
  \item Quiz de teoría (en Intu).
\end{itemize}

\item \textbf{Aplicación en clase (aplicación guiada)}
\begin{itemize}
  \item Cómo se hace el EDA con R: taller guiado corriendo código sobre un dataset \textbf{con problemas intencionales}.
  \begin{itemize}
    \item Diagnóstico: tipos, NA, duplicados, categorías raras, rangos imposibles.
    \item Limpieza: estandarizar nombres/tipos; tratar NA; tratar outliers de forma explícita.
    \item EDA mínima: 1 tabla resumen + 3 gráficos clave.
  \end{itemize}
  \item Producto mínimo: script reproducible de limpieza (con comentarios), dataset ``analysis-ready'' guardado, log breve de decisiones (qué se cambió y por qué), 3 hallazgos + 3 preguntas abiertas.
\end{itemize}
\end{itemize}

%%=======================================%%
%%=======================================%%
\newpage
\section*{Semana 07 - Examen Parcial}
\addcontentsline{toc}{section}{Semana 07 - Examen Parcial}

%%=======================================%%
%%=======================================%%
\newpage
\section{IA aplicada al análisis de datos}

%%=======================================%%
\subsection{Semana 08 - Claude Code / Cursor / VS Code}
\vspace{0.2 cm}

% Nota: la sesión de infraestructura en la nube (creación de VM en AWS/Azure) del
% semestre anterior se elimina del programa (acuerdo de la reunión de ajustes).

\begin{itemize}
\item \textbf{Alcance}
\begin{itemize}
  \item Introducir asistentes de IA para el análisis de datos y la programación: \textbf{Claude Code}, \textbf{Cursor} y \textbf{VS Code} (con asistentes integrados).
  \item Objetivo central: aprender a hacer con IA lo que ya se hizo ``a mano'' en las semanas anteriores (manipulación, limpieza, EDA y visualización).
  \item Conexión con la Semana 02: escribir buenos prompts, ahora en el contexto de código y análisis de datos.
  \item Uso responsable: la IA acelera el trabajo, pero el estudiante debe entender y validar cada paso (conecta con la política de IA del curso).
\end{itemize}

\item \textbf{Fundamentación Teórica (podcast + quiz)}
\begin{itemize}
  \item Podcast: qué son estas herramientas, cómo funcionan (alto nivel), flujos de trabajo típicos, alcances y límites. % TODO: enlace al podcast en Intu
  \item Material para la sesión: % TODO: pegar aquí el/los enlace(s) de podcast/video de esta semana
  \item Quiz de teoría (en Intu).
\end{itemize}

\item \textbf{Aplicación en clase (aplicación guiada)}
\begin{itemize}
  \item Replicar con apoyo de IA un flujo ya conocido: limpieza + EDA de un dataset (diagnóstico, decisiones documentadas, tabla resumen y gráficos).
  \item Buenas prácticas: iterar los prompts, revisar el código generado y verificar los resultados contra lo esperado.
\end{itemize}
\end{itemize}

%%=======================================%%
%%=======================================%%
\newpage
\section*{Semana 09 - Presentación Avance del Proyecto (EDA)}
\addcontentsline{toc}{section}{Semana 09 - Presentación Avance del Proyecto (EDA)}
\vspace{0.2 cm}

\begin{itemize}
  \item Entrega 2 del proyecto final: \textbf{presentación del análisis exploratorio (EDA)}.
  \item Contenido: dada la pregunta de negocio enviada (Entrega 1), cómo se va a responder y qué muestra el EDA de los datos.
  \item Espacio de retroalimentación del profesor y los monitores de cara a la entrega final.
  \item Peso: 5\% de la nota total. Quien no presenta obtiene cero.
  \item Contenido mínimo de la presentación (número de slides, secciones) y rúbrica de evaluación: documento de lineamientos del proyecto.
\end{itemize}

%%=======================================%%
%%=======================================%%
\newpage
\section{Fundamentos de Machine Learning}

%%=======================================%%
\subsection{Semana 10 - Fundamentos de Machine Learning}
\vspace{0.2 cm}

\begin{itemize}
\item \textbf{Alcance}
\begin{itemize}
  \item Entender el objetivo de Machine Learning como \textbf{generalización}: aprender patrones en datos para predecir en datos nuevos.
  \item Construir el \textbf{pipeline} estándar de ML: definir pregunta analítica y variable objetivo (\textit{target}), seleccionar variables predictoras (\textit{features}), separar datos en entrenamiento y prueba (\textit{train/test}), entrenar modelo(s) y ajustar hiperparámetros, evaluar con métricas y comparar contra un baseline, reportar resultados y limitaciones.
  \item Introducir \textbf{validación cruzada} como herramienta para estimar desempeño y reducir sobreajuste.
  \item Identificar errores comunes: \textbf{data leakage}, sobreajuste (\textit{overfitting}) y mala elección de métricas.
\end{itemize}

\item \textbf{Fundamentación Teórica (podcast + videos complementarios + quiz)}
\begin{itemize}
  \item Podcast de la clase. % TODO: enlace al podcast en Intu
  \begin{itemize}
    \item Entrenamiento vs prueba: por qué no se evalúa con los mismos datos con los que se entrena.
    \item Sobreajuste (\textit{overfitting}) vs subajuste (\textit{underfitting}); baseline como el mínimo razonable para saber si el modelo aporta.
    \item Separación de datos: holdout (train/test), principios de muestreo y consideración especial con datos temporales (no mezclar pasado y futuro).
    \item Validación cruzada: idea de $k$-fold CV y cómo se usa para comparar modelos e hiperparámetros.
    \item Riesgos y buenas prácticas: data leakage; el preprocesamiento que se aprende del dataset (escalamiento, imputación) se ajusta \textbf{solo} con train y se aplica a test.
  \end{itemize}
  \item Videos complementarios: construcción de la \textbf{matriz de confusión} (métricas de clasificación) y de \textbf{MAE/RMSE} (métricas de regresión). % Antes ubicados al final de las semanas de clasificación y regresión; se adelantan aquí para que las métricas queden fundamentadas antes de usarlas.
  \item Material para la sesión: % TODO: pegar aquí el/los enlace(s) de podcast/video de esta semana
  \item Quiz de teoría (en Intu).
\end{itemize}

\item \textbf{Aplicación en clase (aplicación guiada)}
\begin{itemize}
  \item Taller guiado con un dataset sencillo (limpio o semi-limpio):
  \begin{itemize}
    \item Definir \textit{target} y \textit{features}; crear partición train/test (por ejemplo 80/20).
    \item Usar predicciones de modelos base previamente entrenados para calcular y comparar métricas de regresión y clasificación.
  \end{itemize}
  \item Producto mínimo: reporte corto con métrica en test, comparación vs baseline, y 2 riesgos/limitaciones.
\end{itemize}
\end{itemize}

%%=======================================%%
%%=======================================%%
\newpage
\section{Aprendizaje supervisado}

%%=======================================%%
\subsection{Semana 11 - Clasificación: pipeline, métricas y modelos base}
\vspace{0.2 cm}

\begin{itemize}
\item \textbf{Alcance}
\begin{itemize}
  \item Definir problemas de \textbf{clasificación} (target categórico) y su pipeline completo.
  \item Métricas de clasificación y decisiones por umbral: matriz de confusión, accuracy (limitaciones), precision/recall y trade-off, (opcional) ROC-AUC como medida agregada.
  \item Modelos a implementar: \textbf{CART} (árbol de clasificación) y \textbf{Random Forest} (ensamble), trabajados en la misma clase.
  \item Buenas prácticas: train/test + validación cruzada para comparar modelos e hiperparámetros.
\end{itemize}

\item \textbf{Fundamentación Teórica (video + quiz)}
\begin{itemize}
  \item Video de la clase. % TODO: enlace al video en Intu
  \begin{itemize}
    \item Pipeline aplicado a clasificación: partición train/test, CV para selección de hiperparámetros, baseline (clase mayoritaria) como referencia.
    \item CART: idea de particiones, profundidad y poda; interpretabilidad vs sobreajuste.
    \item Random Forest: idea de ensamble, \textit{bagging}, aleatoriedad en variables; por qué suele mejorar generalización.
    \item Consideraciones prácticas: desbalance de clases (intuición) y por qué accuracy engaña.
  \end{itemize}
  \item Material para la sesión: % TODO: pegar aquí el/los enlace(s) de podcast/video de esta semana
  \item Quiz de teoría (en Intu).
\end{itemize}

\item \textbf{Aplicación en clase (aplicación guiada)}
\begin{itemize}
  \item Taller:
  \begin{itemize}
    \item Definir target y features; crear split train/test.
    \item Baseline: clase mayoritaria (y métrica correspondiente).
    \item Entrenar CART y evaluar en CV; luego evaluar en test.
    \item Entrenar Random Forest (hiperparámetros básicos) y comparar vs CART.
    \item Reportar métricas y justificar cuál modelo elegir (según costo de errores).
  \end{itemize}
  \item Entregable: reporte corto con comparación (baseline vs CART vs RF) + matriz de confusión y recomendación.
\end{itemize}
\end{itemize}

%%=======================================%%
\newpage
\subsection{Semana 12 - Regresión: pipeline, métricas y modelos base}
\vspace{0.2 cm}

\begin{itemize}
\item \textbf{Alcance}
\begin{itemize}
  \item Problemas de \textbf{regresión} (target continuo) y pipeline completo.
  \item Métricas de regresión: MAE y RMSE (interpretación), comparación contra baseline (media/mediana).
  \item Modelos a implementar (enfoque práctico y comparativo): \textbf{regresión lineal regularizada} (Lasso; Ridge opcional), \textbf{árbol de regresión (CART)}, \textbf{Random Forest} o \textbf{XGBoost} (uno de los dos, según tiempo).
  \item Objetivo pedagógico: entender el trade-off entre interpretabilidad, desempeño y riesgo de sobreajuste.
\end{itemize}

\item \textbf{Fundamentación Teórica (video + quiz)}
\begin{itemize}
  \item Video de la clase. % TODO: enlace al video en Intu
  \begin{itemize}
    \item Pipeline aplicado a regresión: train/test + CV para selección (p.\ ej., lambda en Lasso), baseline (media/mediana) como referencia.
    \item Lasso (y Ridge opcional): idea de regularización y por qué ayuda a generalizar; selección de variables (intuición en Lasso).
    \item CART de regresión y ensambles: árboles (no linealidades e interacciones); Random Forest/XGBoost y por qué mejoran desempeño (visión general).
    \item Interpretación: error en test, diagnóstico simple (residuales), y relevancia de variables (alto nivel).
  \end{itemize}
  \item Material para la sesión: % TODO: pegar aquí el/los enlace(s) de podcast/video de esta semana
  \item Quiz de teoría (en Intu).
\end{itemize}

\item \textbf{Aplicación en clase (aplicación guiada)}
\begin{itemize}
  \item Taller:
  \begin{itemize}
    \item Definir target continuo y features; split train/test.
    \item Baseline: media/mediana y cálculo de MAE/RMSE.
    \item Entrenar Lasso (CV para lambda) y evaluar en test.
    \item Entrenar CART de regresión y comparar.
    \item Entrenar Random Forest o XGBoost y comparar desempeño vs modelos anteriores.
  \end{itemize}
  \item Entregable: tabla comparativa de MAE/RMSE (baseline vs modelos) + recomendación final (con trade-offs).
\end{itemize}
\end{itemize}

%%=======================================%%
%%=======================================%%
\newpage
\section{Aprendizaje no supervisado}

%%=======================================%%
\subsection{Semana 13 - Clustering: fundamentos y métricas}
\vspace{0.2 cm}

% Nota: clustering se mueve al final del bloque de ML (antes iba entre fundamentos
% y clasificación) para que las métricas de la Semana 10 se usen de inmediato en
% las semanas de aprendizaje supervisado (acuerdo de la reunión de ajustes).

\begin{itemize}
\item \textbf{Alcance}
\begin{itemize}
  \item Introducir clustering como herramienta para \textbf{segmentación} (sin target).
  \item Aplicar \textbf{k-means} como algoritmo principal.
  \item Conocer alternativas (alto nivel) y cuándo usarlas: clustering jerárquico (dendrograma), (opcional) DBSCAN para clusters de forma irregular/outliers.
  \item Métricas y selección del número de clusters: elbow / within-cluster sum of squares, silhouette, estabilidad e interpretabilidad.
  \item Producto: segmentos + \textbf{perfilamiento} (tablas y gráficos) para decisiones de negocio.
\end{itemize}

\item \textbf{Fundamentación Teórica (video + quiz)}
\begin{itemize}
  \item Video de la clase (el componente visual es clave para entender cómo se forman los clusters). % TODO: enlace al video en Intu
  \begin{itemize}
    \item Intuición: distancia/similitud, y por qué el escalamiento importa.
    \item k-means: idea de centroides, sensibilidad a escala e inicialización, qué optimiza (compactación intra-cluster).
    \item Elección de $k$: elbow y silhouette (interpretación).
    \item Interpretación: un cluster solo sirve si puede describirse y accionarse.
  \end{itemize}
  \item Material para la sesión: % TODO: pegar aquí el/los enlace(s) de podcast/video de esta semana
  \item Quiz de teoría (en Intu).
\end{itemize}

\item \textbf{Aplicación en clase (aplicación guiada)}
\begin{itemize}
  \item Taller:
  \begin{itemize}
    \item Preparar variables numéricas (selección + escalamiento).
    \item Estimar k-means para varios $k$.
    \item Elegir $k$ con elbow/silhouette (y justificar).
    \item Asignar etiquetas de cluster y perfilar: promedios/medianas por cluster, tamaño de cluster, 2 visualizaciones (p.\ ej., scatter + boxplots por cluster).
  \end{itemize}
  \item Entregable: tabla de perfilamiento + 2 gráficos + 3 insights accionables por segmento.
\end{itemize}
\end{itemize}

%%=======================================%%
%%=======================================%%
\newpage
\section*{Semana 14 - Examen Parcial}
\addcontentsline{toc}{section}{Semana 14 - Examen Parcial}

%%=======================================%%
%%=======================================%%
\newpage
\section*{Semana 15 - Simulacro de la presentación final}
\addcontentsline{toc}{section}{Semana 15 - Simulacro de la presentación final}
\vspace{0.2 cm}

\begin{itemize}
  \item Cada grupo presenta una versión preliminar de su sustentación del proyecto final.
  \item Espacio de retroalimentación de fondo y forma (profesor y monitores) antes de la presentación definitiva.
  \item Objetivo: llegar a la Semana 16 con la presentación ajustada.
\end{itemize}

%%=======================================%%
%%=======================================%%
\newpage
\section*{Semana 16 - Presentación del Proyecto Final}
\addcontentsline{toc}{section}{Semana 16 - Presentación del Proyecto Final}
\vspace{0.2 cm}

\begin{itemize}
  \item Sustentación oral del proyecto final (Entrega 3). Ya no hay entrega de documento escrito: la evaluación final es la presentación.
  \item Equivale a un examen parcial: la asistencia es \textbf{obligatoria} y la ausencia solo se acepta mediante supletorio autorizado por la universidad.
  \item El formato de la presentación es libre; los criterios de evaluación están en la rúbrica del documento de lineamientos del proyecto.
\end{itemize}

\end{document}